%==================================================================
%  SHARED PREAMBLE  --  layout only, no content.
%  Do not edit per job. Edit content/ or a body.tex instead.
%==================================================================
\documentclass[a4paper,10pt]{article}

% \cvroot points at the repo root. Each main.tex sets it before
% loading template/cv.tex; ../../ is the default for variants/* and jobs/*.
\providecommand{\cvroot}{../../}

% Side margins are wider than the top and bottom: a CV is read as a
% column of text, and text that runs to the paper edge reads as dense.
\usepackage[top=1.1cm,bottom=1.1cm,left=1.5cm,right=1.5cm]{geometry}
\usepackage[T1]{fontenc}
\usepackage[utf8]{inputenc}
\usepackage{lmodern}
\usepackage{microtype}
\usepackage{enumitem}
\usepackage{titlesec}
\usepackage{ragged2e}
\usepackage{color}
\usepackage[hidelinks]{hyperref}

% Two neutrals keep the page calm: near-black for headings, a printed
% grey for rules and secondary metadata. Body text stays pure black.
\definecolor{cvink}{gray}{0.10}
\definecolor{cvrule}{gray}{0.60}
\definecolor{cvmuted}{gray}{0.35}

% Machine-readable text layer: helps ATS parsers and copy-paste.
\ifdefined\pdfgentounicode
  \input{glyphtounicode}
  \pdfgentounicode=1
\fi

\urlstyle{same}
\pagestyle{empty}
% Let the last page end where the content ends. With \flushbottom, TeX
% stretches every bit of vertical glue to fill the page, which is what turns
% a nearly-full CV into one with huge, uneven gaps between sections.
\raggedbottom
\setlength{\parindent}{0pt}
\setlength{\parskip}{0pt}
\setlength{\emergencystretch}{3em}
\RaggedRight
% No hyphenation: CV text should never break mid-word, and ATS parsers
% extract "Telecom-mu-ni-ca-tions" as garbage keywords.
\hyphenpenalty=10000
\exhyphenpenalty=10000

%---------------------------- sections ----------------------------
% Sans headings against a serif body give the page a clear hierarchy:
% a recruiter skimming sees the section bands first, then the prose.
% Spacing is fixed (no plus/minus glue) so a sparse CV never opens gaps.
% The gap after the rule is built here instead of in \titlespacing's "after"
% slot. A trailing \penalty10000 immediately before the glue makes that glue
% an illegal page breakpoint: a heading can then never sit alone at the foot
% of a page, which \nopagebreak alone could not prevent (it lands after the
% glue, too late for TeX's page builder).
\titleformat{\section}[block]{\large\sffamily\bfseries\color{cvink}}{}{0pt}{}%
  [\vspace{3pt}{\color{cvrule}\titlerule[0.5pt]}\vspace{5pt}\penalty10000]
\titleformat{name=\section,numberless}[block]{\large\sffamily\bfseries\color{cvink}}{}{0pt}{}%
  [\vspace{3pt}{\color{cvrule}\titlerule[0.5pt]}\vspace{5pt}\penalty10000]
\titlespacing*{\section}{0pt}{11pt}{0pt}
\titlespacing*{name=\section,numberless}{0pt}{11pt}{0pt}

%----------------------------- lists ------------------------------
% Body bullets: a small raised dot, not the full-size class bullet.
\setlist[itemize]{leftmargin=1.15em, labelsep=0.5em,
                  label={\raisebox{0.2ex}{\scriptsize$\bullet$}},
                  topsep=1pt, parsep=0pt, partopsep=0pt, itemsep=2pt}

% Entry list: jobs, degrees and projects. No marker -- a bullet in front
% of every job title is noise, and the bold title already anchors the eye.
\newlist{entrylist}{itemize}{1}
\setlist[entrylist]{label={}, leftmargin=0pt, labelsep=0pt,
                    topsep=0pt, parsep=0pt, partopsep=0pt, itemsep=6pt}

% Achievement bullets nested under an entry. Tight leading, one line of
% separation is enough because the parent spacing already sets the break.
% \interlinepenalty keeps a bullet on one page: a two-line achievement
% that breaks after its first line reads as a mistake in print.
\newlist{subitems}{itemize}{1}
\setlist[subitems]{label={\raisebox{0.2ex}{\scriptsize$\bullet$}},
                   leftmargin=1.15em, labelsep=0.5em,
                   topsep=1pt, parsep=0pt, partopsep=0pt, itemsep=1pt,
                   before={\interlinepenalty=10000}}

%============================ macros ==============================
% Include a fragment from content/.  Usage: \content{experience/minas}
\newcommand{\content}[1]{\input{\cvroot content/#1.tex}}

% Section heading with fixed, honest spacing. Content flows top to bottom and
% spills onto a second page when it needs to; sections are never pushed apart to
% fill leftover space. Use \cvsection instead of \section* everywhere in body.tex.
%
% A negative penalty makes the space before the heading a preferred breakpoint, so a
% heading that would land alone at the foot of a page moves to the next one. It adds
% no glue, so it cannot create a gap.
\newcommand{\cvsection}[1]{\par\penalty-100 \section*{#1}\nopagebreak[3]}

% Page header, built by the renderer:
%   \cvheader{Name}{Headline or empty}{Contact line or empty}{Links line or empty}
% Name, contact and links are centred; a hairline rule closes the block and
% separates it from the first section without spending a whole section's gap.
\newcommand{\cvheader}[4]{%
  \begin{center}%
    {\LARGE\bfseries #1}%
    \if\relax\detokenize{#2}\relax\else\\[2pt]{\large #2}\fi%
    \if\relax\detokenize{#3}\relax\else\\[6pt]{\small #3}\fi%
    \if\relax\detokenize{#4}\relax\else\\[1pt]{\small #4}\fi%
  \end{center}%
  \vspace{-3pt}%
  \noindent{\color{cvrule}\rule{\linewidth}{0.5pt}}%
  \par\vspace{1pt}%
}

% One job entry.  Usage inside \begin{itemize}:
%   \begin{cvjob}{Company}{Location}{Role title}{Dates} ... \end{cvjob}
\newenvironment{cvjob}[4]{%
  \item \textbf{#1} \hfill #2 \\
        \textit{#3} \hfill #4%
  \begin{subitems}%
}{%
  \end{subitems}%
}

% One generic entry for any section (job, degree, award, publication...).
% Usage inside \begin{entrylist}:
%   \cventry{Heading}{Subheading}{Location}{Dates}
% The heading carries the weight; location and dates sit in the muted grey
% so the eye reads title first, then metadata. Follow it with a nested
% \begin{subitems} ... \end{subitems} only when there are bullets; an empty
% nested list is a LaTeX error.
\newcommand{\cventry}[4]{%
  \item {\bfseries\color{cvink}#1}\hfill{\small\textcolor{cvmuted}{#4}}%
  \if\relax\detokenize{#2#3}\relax\else
    \\[1pt]\textit{#2}\hfill{\small\textcolor{cvmuted}{#3}}%
  \fi
}

% One project entry.  Usage inside \begin{itemize}:
%   \begin{cvproject}{Name}{Stack, comma, separated} ... \end{cvproject}
\newenvironment{cvproject}[2]{%
  \item \textbf{#1} \quad {\small #2}%
  \begin{subitems}%
}{%
  \end{subitems}%
}

% One skills line. Each call is its own one-row table, so every label sits on
% a shared baseline and every value is flush at the same x-position without
% requiring a surrounding environment (content/skills/*.tex call it directly).
%   \cvskill{Languages}{C++, C, Python}
\newlength{\cvskilllabel}
\setlength{\cvskilllabel}{8em}
\newcommand{\cvskill}[2]{%
  \par\noindent
  \begin{tabular}{@{}p{\cvskilllabel}@{\hspace{0.7em}}p{\dimexpr\linewidth-\cvskilllabel-0.7em\relax}@{}}%
    \leavevmode\RaggedRight\bfseries\color{cvink}#1: &
    \leavevmode\RaggedRight #2%
  \end{tabular}%
  \par\vspace{2.5pt}%
}

% One education line: \cvedu{Degree}{Institution --- dates}
\newcommand{\cvedu}[2]{\textbf{#1} \hfill #2\par\vspace{2pt}}
